\unitopen{06}
\begin{candos}
\item opowiedzieć o rodzinie: kto jest kim, ile masz rodzeństwa, czy ktoś jest żonaty;
\item użyć przymiotników dzierżawczych (\de{mein, dein, sein, ihr, unser, euer, Ihr}) w mianowniku i bierniku;
\item podać zawód własny i cudzy, także w formie żeńskiej;
\item zapytać o rodzinę i pracę rozmówcy, formalnie i nieformalnie.
\end{candos}

\section{Słownictwo}
\vocabtable{06}

\section{Przymiotniki dzierżawcze}
Odmieniają się \emph{dokładnie jak} \de{ein/kein}: \de{mein} (m, n), \de{meine} (f, mnoga), w bierniku męskim
\de{meinen}.

\begin{center}\small
\begin{tabular}{@{}lll@{}}
\toprule
\textsf{\textbf{Osoba}} & \textsf{\textbf{Przymiotnik}} & \textsf{\textbf{Przykład (mianownik)}}\\
\midrule
\de{ich} & \de{mein} & \de{mein Bruder}, \de{meine Schwester}\\
\de{du} & \de{dein} & \de{dein Vater}, \de{deine Mutter}\\
\de{er, es} & \de{sein} & \de{sein Sohn}, \de{seine Tochter}\\
\de{sie} (ona) & \de{ihr} & \de{ihr Mann}, \de{ihre Kinder}\\
\de{wir} & \de{unser} & \de{unser Team}, \de{unsere Firma}\\
\de{ihr} & \de{euer} & \de{euer Chef}, \de{eure Idee}\\
\de{sie} (oni), \de{Sie} & \de{ihr}, \de{Ihr} & \de{ihr Büro}, \de{Ihr Name}\\
\bottomrule
\end{tabular}
\end{center}
Biernik męski: \de{Ich besuche \textbf{meinen} Bruder.} Żeński i nijaki bez zmian: \de{Ich besuche meine Schwester.}
\de{Ich besuche mein Kind.} Przed \de{euer} końcówka zjada \emph{e}: \de{eure}, \de{euren}.

\begin{falle}
Polskie „swój” nie ma niemieckiego odpowiednika. Mówisz \de{Ich liebe \textbf{meine} Familie}, a o kimś innym
\de{Er liebt \textbf{seine} Familie}: przymiotnik zależy od \emph{właściciela}, a końcówka od rodzaju
\emph{rzeczy posiadanej}. Typowy błąd to \de{Ich besuche mein Bruder}: brat jest męski, a po \de{besuchen} potrzebny
jest biernik: \de{meinen Bruder}.
\end{falle}

\section{Zawody i forma żeńska}
Zawód po \de{sein} nie ma rodzajnika: \de{Ich bin Arzt.} Forma żeńska to zwykle dodanie \de{-in}: \de{der Lehrer
$\rightarrow$ die Lehrerin}; czasem z przegłosem: \de{der Arzt $\rightarrow$ die Ärztin}. Liczba mnoga formy
żeńskiej: \de{-innen} (\de{die Lehrerinnen}). Pytania: \de{Was sind Sie von Beruf?} \de{Was machst du beruflich?}
\de{Als was arbeiten Sie?} Odpowiedź: \de{Ich arbeite als Softwareentwickler.}

\begin{szwajcaria}
W szwajcarskich ogłoszeniach i na wizytówkach zobaczysz \de{Entwickler/in}, \de{Entwickler*in} albo
\de{Entwickler:in}. W rozmowie mówisz po prostu \de{Entwickler} lub \de{Entwicklerin}. Wiele stanowisk zostaje po angielsku
(\de{Senior Backend Engineer}, \de{Product Owner}); przed nazwą stanowiska piszesz rodzajnik, który pasuje do
\emph{osoby}: \de{der Product Owner}, \de{die Product Ownerin}.
\end{szwajcaria}

\begin{praca}
\de{Ich bin Softwareentwickler im Team Backend.} \de{Meine Kollegin ist Architektin.} \de{Unser Teamleiter heißt Anna.}
\de{Mein Chef sitzt in Zürich.} Przy pierwszym spotkaniu pytasz: \de{Wie heißt Ihr Team?} \de{Wie viele Personen arbeiten
in Ihrem Team?} (drugie pytanie za dwa tygodnie).
\end{praca}

\section{Dialogi}
\textbf{Dialog 1 --- nowy kolega}
\begin{dialog}
\turn{A}{Hast du Geschwister?}{Masz rodzeństwo?}
\turn{B}{Ja, ich habe einen Bruder und eine Schwester. Mein Bruder ist Arzt, meine Schwester ist Lehrerin.}{Tak, mam brata i siostrę. Mój brat jest lekarzem, moja siostra nauczycielką.}
\turn{A}{Und deine Eltern? Wo wohnen sie?}{A twoi rodzice? Gdzie mieszkają?}
\turn{B}{Meine Eltern wohnen in Polen. Mein Vater ist Ingenieur, meine Mutter arbeitet als Verkäuferin.}{Moi rodzice mieszkają w Polsce. Mój ojciec jest inżynierem, matka pracuje jako sprzedawczyni.}
\turn{A}{Bist du verheiratet?}{Jesteś żonaty?}
\turn{B}{Nein, ich bin ledig. Aber ich habe eine Freundin.}{Nie, jestem stanu wolnego. Ale mam dziewczynę.}
\end{dialog}

\section{Ćwiczenia}
\begin{cwiczenia}
\exitem{Uzupełnij: \de{Das ist \dots{} Bruder.} (ich)}{\de{mein}}
\exitem{Uzupełnij: \de{Das ist \dots{} Schwester.} (du)}{\de{deine}}
\exitem{Uzupełnij: \de{Ich besuche \dots{} Eltern.} (ich)}{\de{meine}}
\exitem{Uzupełnij: \de{Wir besuchen \dots{} Großvater.} (wir)}{\de{unseren}}
\exitem{Uzupełnij: \de{Er sucht \dots{} Schlüssel.} (er)}{\de{seinen}}
\exitem{Uzupełnij: \de{Sie sucht \dots{} Laptop.} (sie, ona)}{\de{ihren}}
\exitem{Utwórz formę żeńską: \de{der Lehrer}.}{\de{die Lehrerin}}
\exitem{Utwórz formę żeńską: \de{der Arzt}.}{\de{die Ärztin}}
\exitem{Utwórz liczbę mnogą formy żeńskiej: \de{die Studentin}.}{\de{die Studentinnen}}
\exitem{Przetłumacz: „Mam dwóch braci i jedną siostrę.”}{\de{Ich habe zwei Brüder und eine Schwester.}}
\exitem{Przetłumacz: „Moja siostra jest lekarką.”}{\de{Meine Schwester ist Ärztin.}}
\exitem{Zapytaj kolegę, czym zajmuje się jego ojciec.}{\de{Was ist dein Vater von Beruf?}}
\exitem{Przetłumacz: „Pracuję jako programista.”}{\de{Ich arbeite als Softwareentwickler.}}
\end{cwiczenia}

\section{Trening słownictwa (generowany)}
\drill{06}

\weekend{Przygotuj 60 sekund o rodzinie (kto, ile lat, czym się zajmuje) i 60 sekund o zespole w pracy
(\de{Mein Kollege \dots{} ist \dots{}}). Nagraj i wyślij do korekty. Zwróć uwagę na biernik męski po czasownikach.}
